\section{Conclusion}

Thermal sensing for nighttime segmentation is usually treated as an all-or-nothing
input: either the full thermal stream is fused into the network, or its knowledge is
distilled into RGB. We have shown that neither is necessary. On MFNet, the nighttime
pedestrian advantage of a full-resolution thermal teacher is modality-borne and does
not distil into an RGB student, but one sixteenth of the thermal pixels recovers 81\%
of that advantage. A globally fused low-resolution model still loses overall, because
thermal is worse on the other eight classes. A small person-aware gate resolves the
conflict: it routes thermal evidence only where the experts disagree and a pedestrian
is at stake, and it raises nighttime pedestrian IoU from 45.24 to $61.36\pm0.39$ while
improving overall mIoU from 56.38 to $57.57\pm0.08$. Under corrupted RGB the full gate
is the best overall model, and the thermal expert remains useful when RGB collapses.

The broader lesson is that sensor budgets should be measured, not assumed. The
$120\times160$ knee is a property of MFNet's pedestrian scale and capture conditions,
not a universal constant; a PST900 pilot shows that smaller, lower-contrast survivors
need more thermal pixels. The right deployment procedure is therefore to sweep the
budget on the target domain, select the routing constraint from the safety class, and
report the compute cost alongside the sensor cost. We hope this work makes that
procedure the default for RGB-T perception on low-cost platforms.
